
\section*{\SGAInlineRefTarget{sga3:ii:section:3}3. The tangent bundle, condition (E)}
\addcontentsline{toc}{section}{3. The tangent bundle, condition (E)}

Throughout this section, unless otherwise stated, the letters
\(\mathcal M,\mathcal M',\mathcal N,\ldots\) will denote finite free
\(\mathcal O_S\)-modules (that is, modules isomorphic to finite direct sums
of copies of \(\mathcal O_S\)).

We shall systematically use the identifications established in Expos\'e~I.
Thus, by a ``functor over \(S\)'' we shall mean indifferently either a
functor equipped with a morphism to \(S=h_S\), or a functor on the category
of objects over \(S\). We shall use ``group functor over \(S\),'' and
similar expressions, in the same way.

\SGAInlineRefTarget{sga3:ii:definition:3-1}\subsection*{\SGARefLink{sga3:ii:definition:3-1}{Definition 3.1}.}
Let \(S\) be a scheme, let \(\mathcal M\) be a finite free
\(\mathcal O_S\)-module, and let \(X\) be a functor over \(S\). The
\emph{tangent bundle of \(X\) over \(S\), relative to the
\(\mathcal O_S\)-module \(\mathcal M\)}, is the \(S\)-functor
\[
  T_{X/S}(\mathcal M)
  =\underline{\operatorname{Hom}}_S(I_S(\mathcal M),X).
\]
In particular, the \emph{tangent bundle of \(X\) over \(S\)} is
\[
  T_{X/S}=T_{X/S}(\mathcal O_S)
  =\underline{\operatorname{Hom}}_S(I_S,X).
\]

The assignment \(\mathcal M\mapsto T_{X/S}(\mathcal M)\) is a covariant
functor from finite free \(\mathcal O_S\)-modules to \(S\)-functors. In
particular, the morphisms \(\mathcal M\to0\) and \(0\to\mathcal M\)
define, respectively, an \(S\)-morphism
\[
  T_{X/S}(\mathcal M)\longrightarrow
  T_{X/S}(0)\simeq X
\]
and a section of it, called the \emph{zero section}.\footnote{Translator's
note: the 1970 text has ``\(0\to\mathcal M\) and \(\mathcal M\to0\)'' in
this order, as for \(I_S(\mathcal M)\) in \S2. Since
\(\mathcal M\mapsto T_{X/S}(\mathcal M)\) is covariant, the projection comes
from \(\mathcal M\to0\) and the section from \(0\to\mathcal M\); the
Gille--Polo re-edition (2011), Expos\'e~II, corrects the order.}

It follows from this that \(\mathcal M\mapsto T_{X/S}(\mathcal M)\) is a
covariant functor from finite free \(\mathcal O_S\)-modules to functors
over \(X\). In particular, \(\mathcal O(S)\) is a set of operators of the
\(X\)-functor \(T_{X/S}(\mathcal M)\), compatible with functoriality in
\(\mathcal M\).

\SGAInlineRefTarget{sga3:ii:definition:3-2}\subsection*{\SGARefLink{sga3:ii:definition:3-2}{Definition 3.2}.}
Let \(u\in X(S)=\operatorname{Hom}_S(S,X)=\Gamma(X/S)\). The
\emph{tangent space of \(X\) over \(S\) at \(u\), relative to
\(\mathcal M\)}, denoted \(L^u_{X/S}(\mathcal M)\), is the \(S\)-functor
obtained from the \(X\)-functor \(T_{X/S}(\mathcal M)\) by pullback along
\(u:S\to X\).

\SGAThreeNativeDiagram[0.36\textwidth]
  {native_diagrams/exp2/exp2_p008_tangent_space_cartesian_square.tex}
  {The cartesian square defining \(L^u_{X/S}(\mathcal M)\).}
  {fig:exp2-tangent-space-cartesian}
  {LNM 151, pp.~48--50}

In particular, \(L^u_{X/S}(\mathcal O_S)\) is denoted \(L^u_{X/S}\) and
is called the \emph{tangent space of \(X\) over \(S\) at \(u\)}.\footnote{Translator's
note: the Gille--Polo re-edition (2011), Expos\'e~II, makes the two
structure morphisms explicit (its 3.1.1 and 3.2.1). In our words: a point of
\(T_{X/S}(\mathcal M)\) with values in \(S'\) is an \(S\)-morphism
\(f\colon S'\times_SI_S(\mathcal M)\to X\); its projection in \(X(S')\) is
\(f\circ(\operatorname{id}_{S'}\times\varepsilon_{\mathcal M})\), and the
zero section sends \(g\in X(S')\) to
\(g\circ(\operatorname{id}_{S'}\times\rho)\). Hence, for \(t\colon S'\to S\),
\(L^u_{X/S}(\mathcal M)(S')\) is the set of those \(f\) with
\(f\circ(\operatorname{id}_{S'}\times\varepsilon_{\mathcal M})=u\circ t\).}

\SGAInlineRefTarget{sga3:ii:proposition:3-3}\subsection*{\SGARefLink{sga3:ii:proposition:3-3}{Proposition 3.3}.}
If \(X\) is represented by an \(S\)-scheme \(\widetilde X\), then
\(T_{X/S}(\mathcal M)\) and \(L^u_{X/S}(\mathcal M)\) are representable.
In particular, \(T_{X/S}\) and \(L^u_{X/S}\) are represented by the
vector bundles
\[
  V(\Omega^1_{\widetilde X/S})
  \qquad\text{and}\qquad
  V(u^*\Omega^1_{\widetilde X/S})
\]
over \(\widetilde X\) and \(S\), respectively.

It plainly suffices to prove the assertion for \(T_{X/S}(\mathcal M)\),
since the corresponding assertion for \(L^u_{X/S}(\mathcal M)\) follows
by pullback. By \SGARefLink{sga3:ii:corollary:2-2-1}{2.2, Corollary}, it even suffices to treat \(T_{X/S}\),
and in that case the proposition is none other than the remark made in the
\SGARefLink{sga3:ii:nota:2-2-2}{Nota} after 2.2.

\SGARefTarget{sga3:ii:remark:3-3-1}%
The proposition gives a particularly simple description of the vector
bundle representing \(L^u_{X/S}\). The image of the section \(u\) of
\(X\) over \(S\) is locally closed,\footnote{Translator's note: the
Gille--Polo re-edition (2011), Expos\'e~II, refers here to EGA~I, 5.3.11: the
graph of an \(S\)-morphism is an immersion. The section \(u\colon S\to X\) is
its own graph, hence an immersion.}
and hence is defined by a quasi-coherent ideal \(\mathfrak m\) in the
scheme induced on an open subscheme of \(X\). The quotient
\(\mathfrak m/\mathfrak m^2\) may be regarded as a quasi-coherent module
on \(S\); it is this module that defines the vector bundle in question.

For example, let \(X\) be an algebraic scheme over a field \(k\), and let
\(u\) be a \(k\)-rational point of \(X\). Then
\(L^u_{X/k}=W(\mathfrak t)\), where \(\mathfrak t\) is the Zariski tangent
\(k\)-vector space of the local ring \(\mathcal O_{X,u}\).\footnote{Translator's
note: the 1970 text has \(L^u_{X/k}=V(\mathfrak t)\). With the conventions
of Expos\'e~I (\SGARefLink{sga3:i:subsection:4-6}{I~4.6}), \(V(\mathcal E)\) is the spectrum of the
symmetric algebra of \(\mathcal E\) and its points are the linear forms on
\(\mathcal E\); the representing bundle is \(V(\mathfrak m/\mathfrak m^2)\)
by the preceding sentence, \(\mathfrak m\) being the maximal ideal of
\(\mathcal O_{X,u}\), and since \(\mathfrak t\) is the dual of the
finite-dimensional space \(\mathfrak m/\mathfrak m^2\), this is
\(W(\mathfrak t)\). The Gille--Polo re-edition (2011), Expos\'e~II, corrects
accordingly.}

Returning to the general situation, observe first that
\(L^u_{X/S}(\mathcal M)\) is covariant from finite free
\(\mathcal O_S\)-modules to functors over \(S\). In particular,
\(\mathcal O(S)\) is a set of operators of the \(S\)-functor
\(L^u_{X/S}(\mathcal M)\), compatible with functoriality in
\(\mathcal M\).

\SGAInlineRefTarget{sga3:ii:proposition:3-4}\subsection*{\SGARefLink{sga3:ii:proposition:3-4}{Proposition 3.4}.}
Formation of \(T_{X/S}(\mathcal M)\) and \(L^u_{X/S}(\mathcal M)\)
commutes with base extension: let \(S'\) be a scheme over \(S\); there are
isomorphisms, functorial in \(\mathcal M\),
\[
  T_{X_{S'}/S'}(\mathcal M\otimes_{\mathcal O_S}\mathcal O_{S'})
    \xrightarrow{\sim} T_{X/S}(\mathcal M)_{S'},
\]
\[
  L^{u'}_{X_{S'}/S'}(\mathcal M\otimes_{\mathcal O_S}\mathcal O_{S'})
    \xrightarrow{\sim} L^u_{X/S}(\mathcal M)_{S'},
  \qquad u'=u_{S'}.
\]
This follows immediately from the compatibility of
\(\underline{\operatorname{Hom}}\) with base extension.

\SGAInlineRefTarget{sga3:ii:corollary:3-4-1}\subsection*{Corollary.}
The \(X\)-functor \(T_{X/S}(\mathcal M)\), respectively the \(S\)-functor
\(L^u_{X/S}(\mathcal M)\), is naturally an object with operators
\(\mathcal O_X\), respectively \(\mathcal O_S\), and these structures
are functorial in \(\mathcal M\).

For \(L^u_{X/S}(\mathcal M)\), this is immediate. For each \(S'\) over
\(S\), the ring \(\mathcal O(S')\) acts on
\(\mathcal M\otimes\mathcal O_{S'}\), hence on
\[
  L^{u'}_{X_{S'}/S'}(\mathcal M\otimes\mathcal O_{S'})
  =L^u_{X/S}(\mathcal M)_{S'};
\]
this action is functorial in \(S'\).

The case of \(T_{X/S}(\mathcal M)\) is a little more complicated. We must
equip each \(\bigl(T_{X/S}(\mathcal M)\bigr)_{X'}(X')\) (\(X'\) over \(X\))
with a structure of set with a set of operators \(\mathcal O(X')\),
functorially in \(X'\). For this one constructs the following diagram,
where \(X_{X'}\) denotes \(X\times_SX'\), and \(f'\) the section of
\(X_{X'}\) over \(X'\) defined by \(f:X'\to X\).

\SGAThreeNativeDiagram[0.39\textwidth]
  {native_diagrams/exp2/exp2_p010_base_change_operator_diagram.tex}
  {The base-change diagram used to define the \(\mathcal O_X\)-operator structure.}
  {fig:exp2-base-change-operator}
  {LNM 151, pp.~50--51}

This diagram shows that \(\bigl(T_{X/S}(\mathcal M)\bigr)_{X'}(X')\)
identifies with \(L^{f'}_{X_{X'}/X'}(\mathcal M)(X')\), on which
\(\mathcal O(X')\) operates. It only remains to check that this
construction is functorial in \(X'\), which presents no difficulty.\footnote{Translator's
note: the Gille--Polo re-edition (2011), Expos\'e~II, describes these
operators directly (its 3.4.2 and 3.4.3). In our words: an element of
\(\operatorname{Hom}_X(X',T_{X/S}(\mathcal M))\) is an \(S\)-morphism
\(\varphi\colon I_S(\mathcal M)\times_SX'\to X\) whose composite with the
zero section \(X'\to I_S(\mathcal M)\times_SX'\) is \(f\), and
\(a\in\mathcal O(X')\) sends \(\varphi\) to \(\varphi\circ a^*\), where
\(a^*\) is the endomorphism \((m,x')\mapsto(m\cdot a(x'),x')\) of
\(I_S(\mathcal M)\times_SX'\) (see our note on base change in \S2); since
\(a^*\) preserves the zero section, \(\varphi\circ a^*\) again lies over
\(f\). In the same way \(a\in\mathcal O(S')\) acts on
\(L^u_{X/S}(\mathcal M)(S')\) by composition with \(a^*\).}

The isomorphisms of the preceding proposition are then, by construction,
isomorphisms for the structures of \(\mathcal O_{X_{S'}}\)-objects,
respectively \(\mathcal O_{S'}\)-objects.

\SGAInlineRefTarget{sga3:ii:definition:3-5}\subsection*{\SGARefLink{sga3:ii:definition:3-5}{Definition 3.5}.}
Let \(S\) be a scheme and \(X\) an \(S\)-functor. We say that \(X\)
\emph{satisfies condition (E) relative to \(S\)} if, for every \(S'\)
over \(S\) and every pair of finite free \(\mathcal O_{S'}\)-modules
\(\mathcal M,\mathcal N\), the diagram of sets obtained by applying \(X\)
to diagram \((*)\) of \SGARefLink{sga3:ii:definition:2-1}{2.1} is cartesian.\footnote{Translator's
note: the Gille--Polo re-edition (2011), Expos\'e~II, refers here to
\SGARefLink{sga3:ii:subsection:6-2}{6.2} and \SGARefLink{sga3:ii:subsection:6-3}{6.3} for functors and groups that do not satisfy~(E).}

\SGAThreeNativeDiagram[0.46\textwidth]
  {native_diagrams/exp2/exp2_p011_condition_E_cartesian_diamond.tex}
  {The cartesian diagram in condition (E).}
  {fig:exp2-condition-e}
  {LNM 151, p.~52}

\subsection*{\SGAInlineRefTarget{sga3:ii:subsection:3-5-1}3.5.1.}
Equivalently, the functor
\(\mathcal M\mapsto T_{X/S}(\mathcal M)\) carries direct sums of finite
free \(\mathcal O_S\)-modules to products of \(X\)-functors.
\SGARefLink{sga3:ii:subsection:2-2}{Proposition~2.2} shows that every representable functor satisfies
condition (E). If \(X\) satisfies condition (E) relative to \(S\), the
functor \(\mathcal M\mapsto T_{X/S}(\mathcal M)\), commuting with products,
carries groups to groups. In particular, \(T_{X/S}(\mathcal M)\) is a
commutative \(X\)-group; for the same reason,
\(L^u_{X/S}(\mathcal M)\) is a commutative \(S\)-group.\footnote{Translator's
note: the Gille--Polo re-edition (2011), Expos\'e~II, states the reason: if
\(\mathcal M\mapsto T_{X/S}(\mathcal M)\) carries direct sums to products
over \(X\), so does \(\mathcal M\mapsto
L^u_{X/S}(\mathcal M)=S\times_XT_{X/S}(\mathcal M)\), since base change
along \(u\) commutes with fibre products.}

\SGAInlineRefTarget{sga3:ii:proposition:3-6}\subsection*{\SGARefLink{sga3:ii:proposition:3-6}{Proposition 3.6}.}
If \(X/S\) satisfies (E), then the abelian-group structure on
\(T_{X/S}(\mathcal M)\), respectively on \(L^u_{X/S}(\mathcal M)\),
together with the action of \(\mathcal O_X\), respectively of
\(\mathcal O_S\), makes it an \(\mathcal O_X\)-module, respectively an
\(\mathcal O_S\)-module.

The operator action is functorial in \(\mathcal M\), and therefore
respects the abelian-group structure obtained by functoriality from that
of \(\mathcal M\).\footnote{Translator's note: the Gille--Polo re-edition
(2011), Expos\'e~II, details this proof. In our words: by the construction
of the operators (Corollary of 3.4) it suffices, after base change, to treat
\(L^u_{X/S}(\mathcal M)\) and its points over \(S\), on which
\(\lambda\in\mathcal O(S)\) acts through the homothety \(\lambda\) of
\(\mathcal M\). Writing \(L(\mathcal M)\) for \(L^u_{X/S}(\mathcal M)\), under (E) the two
projections identify \(L(\mathcal M\oplus\mathcal M)\) with
\(L(\mathcal M)\times L(\mathcal M)\);
then the sum \(\sigma\colon\mathcal M\oplus\mathcal M\to\mathcal M\) induces
the addition, the diagonal \(\delta\) the diagonal, and
\(\lambda\oplus\mu\) the map \(L(\lambda)\times L(\mu)\). From
\(\sigma\circ(\lambda\oplus\lambda)=\lambda\circ\sigma\) one gets
\(\lambda(x+y)=\lambda x+\lambda y\), and from
\(\lambda+\mu=\sigma\circ(\lambda\oplus\mu)\circ\delta\) one gets
\((\lambda+\mu)x=\lambda x+\mu x\); the other axioms follow from
functoriality in \(\mathcal M\).}

\SGAInlineRefTarget{sga3:ii:remark:3-6-1}\subsection*{Remark.}
If \(X\) is representable, then it satisfies (E), and \(T_{X/S}\) and
\(L^u_{X/S}\) are represented by vector bundles. The preceding laws
are then exactly those induced by the vector-bundle structures
(cf. \SGARefLink{sga3:i:subsection:4-6}{I~4.6}).\footnote{Translator's note: the Gille--Polo
re-edition (2011), Expos\'e~II, makes this explicit: for \(f\colon X'\to X\),
the bijection
\(\operatorname{Hom}_X(X',T_{X/S}(\mathcal M))\simeq
\operatorname{Hom}_{\mathcal O_{X'}}(f^*\Omega^1_{X/S},
\mathcal M\otimes_{\mathcal O_S}\mathcal O_{X'})\) given by the proof of
SGA~1, III~5.1 carries the operators of \(\mathcal O(X')\) to the obvious
\(\mathcal O(X')\)-module structure of the right-hand side.}

\SGAInlineRefTarget{sga3:ii:proposition:3-4-bis}\subsection*{\SGARefLink{sga3:ii:proposition:3-4-bis}{Proposition 3.4 bis}.}
If \(X\) satisfies condition (E) relative to \(S\), then \(X_{S'}\)
satisfies condition (E) relative to \(S'\), and the isomorphisms of 3.4
respect the structures of \(\mathcal O_{X_{S'}}\)-modules, respectively of
\(\mathcal O_{S'}\)-modules.

No further comment is needed.

Abbreviation: instead of saying ``\(X\) satisfies condition (E) relative
to \(S\)'', we shall sometimes say ``\(X/S\) satisfies condition (E)''.

\SGAInlineRefTarget{sga3:ii:proposition:3-7}\subsection*{\SGARefLink{sga3:ii:proposition:3-7}{Proposition 3.7}.}
The functors \(T_{X/S}(\mathcal M)\) and \(L^u_{X/S}(\mathcal M)\) are
functorial in \(X\): if \(f:X\to X'\) is an \(S\)-morphism, the
following diagrams commute.

\SGAThreeNativeDiagram[0.74\textwidth]
  {native_diagrams/exp2/exp2_p012_functoriality_T_L_diagrams.tex}
  {Functoriality of tangent bundles and tangent spaces in \(X\).}
  {fig:exp2-tangent-functoriality}
  {LNM 151, pp.~52--53}

The morphisms \(T(f)\) and \(L(f)\) follow immediately from the
definitions. Commutativity follows from functoriality
in \(\mathcal M\) and from \(T_{X/S}(0)=X\).\footnote{Translator's note:
the Gille--Polo re-edition (2011), Expos\'e~II, adds that if \(f\) is a
monomorphism, so are \(T(f)\) and \(L(f)\): on points, \(T(f)\) is
composition with \(f\) of morphisms \(I_S(\mathcal M)\times_SS'\to X\),
hence injective when \(f\) is, and \(L(f)\) is induced by \(T(f)\) on the
fibres over \(u\) and \(f\circ u\).}

\SGAInlineRefTarget{sga3:ii:remark:3-7-1}\subsection*{Remark.}
The square above is cartesian when \(f\) is an open immersion, and more
generally when \(f\) is \'etale.\footnote{Translator's note: the
Gille--Polo re-edition (2011), Expos\'e~II, assumes here that \(X\) and
\(X'\) are representable; it then deduces from the Nota of 2.2 that the
square is cartesian as soon as \(f^*\Omega^1_{X'/S}\to\Omega^1_{X/S}\) is an
isomorphism, and that \(L^u_{X/S}(\mathcal M)\simeq
L^{f\circ u}_{X'/S}(\mathcal M)\) in that case. The added hypothesis is not
needed for the statement above, \(f\) being an open immersion or an \'etale
morphism (that is, representable by \'etale morphisms of schemes): a point
of \(T_{X'/S}(\mathcal M)\) with values in \(S'\) is a morphism
\(I_S(\mathcal M)\times_SS'\to X'\); if its restriction to the closed
subscheme \(S'\), defined by an ideal of square zero, factors through \(X\),
then it factors uniquely through \(X\), since \(f\) is formally \'etale
(EGA~IV\textsubscript{4}, 17.1.1 and 17.3.1).} In general, it defines a morphism of
\(X\)-functors:

\SGAThreeNativeDiagram[0.42\textwidth]
  {native_diagrams/exp2/exp2_p013_general_X_functor_triangle.tex}
  {The induced morphism over \(X\) in the general case.}
  {fig:exp2-general-x-functor}
  {LNM 151, p.~53}

\SGAInlineRefTarget{sga3:ii:proposition:3-7-bis}\subsection*{\SGARefLink{sga3:ii:proposition:3-7-bis}{Proposition 3.7 bis}.}
If \(X\) and \(X'\) satisfy
(E) relative to \(S\), then
\[
  T_{X/S}(\mathcal M)\xrightarrow{\ T(f)\ }
  T_{X'/S}(\mathcal M)_X,
  \qquad\text{respectively}\qquad
  L^u_{X/S}(\mathcal M)\xrightarrow{\ L(f)\ }
  L^{f\circ u}_{X'/S}(\mathcal M),
\]
is a morphism of \(\mathcal O_X\)-modules, respectively of
\(\mathcal O_S\)-modules.

This follows from \SGARefLink{sga3:ii:proposition:3-7}{Proposition~3.7} by functoriality in \(\mathcal M\).

\SGAInlineRefTarget{sga3:ii:proposition:3-8}\subsection*{\SGARefLink{sga3:ii:proposition:3-8}{Proposition 3.8}.}
For two functors \(X,Y\) over \(S\), there are isomorphisms, functorial in
\(\mathcal M\),
\[
\begin{aligned}
  T_{X\times_SY/S}(\mathcal M)
    &\xrightarrow{\sim}T_{X/S}(\mathcal M)\times_ST_{Y/S}(\mathcal M),\\
  L^{(u,v)}_{X\times_SY/S}(\mathcal M)
    &\xrightarrow{\sim}L^u_{X/S}(\mathcal M)\times_SL^v_{Y/S}(\mathcal M).
\end{aligned}
\]
This is evident from the definitions.\footnote{Translator's note: the
Gille--Polo re-edition (2011), Expos\'e~II, writes these isomorphisms in the
opposite direction and justifies them: the first is the compatibility of
\(\underline{\operatorname{Hom}}_S(I_S(\mathcal M),-)\) with products (the
last remark of \S1), and the second follows from it by base change along
\((u,v)\colon S\to X\times_SY\).}

\SGARefTarget{sga3:ii:remark:3-8-0}%
Let us observe that the first isomorphism may also be read as an
isomorphism of \(X\times_SY\)-functors
\[
  T_{X\times_SY/S}(\mathcal M)
  \xrightarrow{\sim}
  \bigl(T_{X/S}(\mathcal M)\times_X(X\times_SY)\bigr)
  \mathop{\times}_{X\times_SY}
  \bigl(T_{Y/S}(\mathcal M)\times_Y(X\times_SY)\bigr).
\]

\subsection*{Corollary.}
If \(X/S\) carries an algebraic structure defined by finite cartesian
products, then \(T_{X/S}(\mathcal M)\) carries a structure of the same
kind, and the projection \(T_{X/S}(\mathcal M)\to X\) is a morphism for
that structure.

\SGAInlineRefTarget{sga3:ii:proposition:3-8-bis}\subsection*{\SGARefLink{sga3:ii:proposition:3-8-bis}{Proposition 3.8 bis}.}
If \(X\) and \(Y\) satisfy condition (E) relative to \(S\), then so does
\(X\times_SY\), and the isomorphisms of 3.8 respect the module
structures.\footnote{Translator's note: the Gille--Polo re-edition (2011),
Expos\'e~II, adds a proof. In our words: by 3.8, the functor
\(\mathcal M\mapsto T_{X\times_SY/S}(\mathcal M)\) is the fibre product over
\(S\) of \(\mathcal M\mapsto T_{X/S}(\mathcal M)\) and
\(\mathcal M\mapsto T_{Y/S}(\mathcal M)\); as these carry direct sums to
products over \(X\) and over \(Y\), it carries direct sums to products over
\(X\times_SY\), which is (E) by 3.5.1. For \(Z\) over \(X\times_SY\), the
operators of \(a\in\mathcal O(Z)\) are given on both sides by composition
with the same endomorphism \(a^*\) of \(I_S(\mathcal M)\times_SZ\) (see the
note to the Corollary of 3.4), and a pair \((\varphi,\psi)\) of morphisms to
\(X\) and to \(Y\) satisfies
\((\varphi\circ a^*,\psi\circ a^*)=(\varphi,\psi)\circ a^*\).}

\SGARefTarget{sga3:ii:definition:3-9-0-1}%
To state conveniently the properties that follow, let us introduce the
following terminology: an \emph{\(H\)-set} is a set equipped with a
composition law having a two-sided identity; an \emph{\(H\)-object} in a
category \(\mathscr C\) is defined in the usual way: it is thus an object
\(X\) of \(\mathscr C\) equipped with a morphism \(X\times X\to X\) such
that there exists a section of \(X\) (over the final object) having the
properties of a two-sided identity. Every \(\mathscr C\)-monoid, in
particular every \(\mathscr C\)-group, is thus a
\(\mathscr C\)-\(H\)-object. In particular, an \(H\)-object of the category
of functors over the scheme \(S\) will be called an
\emph{\(S\)-\(H\)-functor}. If \(X\) is an \(S\)-\(H\)-functor, then
\(T_{X/S}(\mathcal M)\) is an \(S\)-\(H\)-functor, and if we write
\[
  \underline{\operatorname{Lie}}(X/S,\mathcal M)=L^e_{X/S}(\mathcal M),
  \qquad
  \underline{\operatorname{Lie}}(X/S)
  =\underline{\operatorname{Lie}}(X/S,\mathcal O_S),
\]
where \(e\) denotes the unit section of \(X\), then
\(\underline{\operatorname{Lie}}(X/S,\mathcal M)\) is also an
\(S\)-\(H\)-functor.\footnote{Translator's note: the Gille--Polo re-edition
(2011), Expos\'e~II, introduces this terminology through the case of an
\(S\)-group and adds a corollary, used in 4.1. In our words: the structure
of \(T_{X/S}(\mathcal M)\) is obtained by applying
\(\underline{\operatorname{Hom}}_S(I_S(\mathcal M),-)\) to that of \(X\)
(the last remark of \S1); the projection \(p\colon T_{X/S}(\mathcal M)\to X\)
and the zero section \(s\), induced by \(\varepsilon_{\mathcal M}\) and by
\(\rho\), are therefore morphisms of \(S\)-\(H\)-functors (of \(S\)-groups if
\(X\) is an \(S\)-group), \(s\) is a section of \(p\), and
\(\underline{\operatorname{Lie}}(X/S,\mathcal M)\) is the kernel of \(p\),
that is, the inverse image of the unit section of \(X\). Proposition 3.9
compares its law with the group law given by~(E).}

\SGAInlineRefTarget{sga3:ii:proposition:3-9}\subsection*{\SGARefLink{sga3:ii:proposition:3-9}{Proposition 3.9}.}
Let \(X\) be an \(S\)-\(H\)-functor satisfying (E) relative to \(S\).
The \(S\)-\(H\)-functor structure on
\(\underline{\operatorname{Lie}}(X/S,\mathcal M)\) induced by that of
\(X\) agrees with the \(S\)-group structure defined in \SGARefLink{sga3:ii:subsection:3-5-1}{3.5.1}.

The preceding discussion shows that
\(\underline{\operatorname{Lie}}(X/S,\mathcal M)\) is an \(H\)-object in
the category of \(\mathcal O_S\)-modules. The proposition follows from
the next lemma.

\SGAInlineRefTarget{sga3:ii:lemma:3-10}\subsection*{\SGARefLink{sga3:ii:lemma:3-10}{Lemma 3.10}.}
Let \(\mathscr C\) be a category. Let \(G\) be an \(H\)-object in the
category of \(\mathscr C\)-\(H\)-objects. Thus \(G\) is a
\(\mathscr C\)-\(H\)-object, whose composition law we denote
\(f:G\times G\to G\), together with a morphism of
\(\mathscr C\)-\(H\)-objects \(h:G\times G\to G\).
Then \(f=h\), and \(f\) is commutative.

Evaluating the functors on a variable argument reduces the proof, as
usual, to the case in which \(\mathscr C\) is the category of sets. We
then have a set \(G\) and two maps \(f,h:G\times G\to G\) satisfying
\[
  h(f(x,y),f(z,t))=f(h(x,z),h(y,t)).
\]
Let \(e,u\in G\) be the respective identities for \(f\) and \(h\), so
that
\[
  f(e,x)=f(x,e)=x,
  \qquad h(u,x)=h(x,u)=x.
\]
One sees first that
\[
  h(f(u,y),f(x,u))=h(f(x,u),f(u,y))=f(x,y).
\]
In particular, for \(y=e\), one obtains
\[
  x=f(x,e)=h(f(u,e),f(x,u))=h(u,f(x,u))=f(x,u),
\]
whence, substituting in the original equality,
\[
  f(x,y)=h(x,y)=h(y,x).\footnotemark
\]
\footnotetext{Translator's note: here \(G\times G\) carries the product law
\(((x,z),(y,t))\mapsto(f(x,y),f(z,t))\), so that \(h\) being a morphism
means \(h(f(x,y),f(z,t))=f(h(x,z),h(y,t))\); and the substitution also uses
\(f(u,y)=y\), obtained in the same way for \(x=e\). The Gille--Polo
re-edition (2011), Expos\'e~II, writes out both points.}

\SGAInlineRefTarget{sga3:ii:corollary:3-9-1}\subsection*{Corollary 1.}
If \(X\) is an \(S\)-\(H\)-functor satisfying (E) relative to \(S\),
every element of \(X(I_S(\mathcal M))\) that maps to the identity element
of \(X(S)\) is invertible.

\SGAInlineRefTarget{sga3:ii:corollary:3-9-2}\subsection*{Corollary 2.}
If \(X\) is an \(S\)-monoid satisfying (E) relative to \(S\), an element
of \(X(I_S(\mathcal M))\) is invertible if and only if its image in
\(X(S)\) is invertible.

\SGAInlineRefTarget{sga3:ii:corollary:3-9-3}\subsection*{Corollary 3.}
If \(X\) is an \(S\)-group satisfying (E) relative to \(S\), the two
\(S\)-group laws on
\(\underline{\operatorname{Lie}}(X/S,\mathcal M)\) coincide.

\par\medskip\noindent
\SGARefTarget{sga3:ii:corollary:3-9-4}%
\emph{Addition} (after the Gille--Polo re-edition (2011), Expos\'e~II; our
formulation). Let \(G\) be an \(S\)-group satisfying (E) relative to \(S\),
let \(n\in\mathbb Z\), and let \(n_G\colon G\to G\) be the morphism of
\(S\)-functors \(g\mapsto g^n\). It is not a group morphism in general, but
it preserves the unit section, so it induces (3.7) a morphism
\(L(n_G)\colon\underline{\operatorname{Lie}}(G/S)\to
\underline{\operatorname{Lie}}(G/S)\); this is multiplication by \(n\).
Indeed, \(T(n_G)\) is the \(n\)-th power map of the \(S\)-group
\(T_{G/S}\), whose law is obtained by applying
\(\underline{\operatorname{Hom}}_S(I_S,-)\) to that of \(G\); so
\(L(n_G)\) is the \(n\)-th power map of
\(\underline{\operatorname{Lie}}(G/S)\) for the law coming from \(G\). By
Corollary~3 this law is the addition of the \(\mathcal O_S\)-module
\(\underline{\operatorname{Lie}}(G/S)\), for which the \(n\)-th power of
\(x\) is \(nx\).
\par\medskip

Before drawing further consequences from the preceding proposition, let
us prove another functoriality result.

\SGAInlineRefTarget{sga3:ii:proposition:3-11}\subsection*{\SGARefLink{sga3:ii:proposition:3-11}{Proposition 3.11}.}
In the situation of \SGARefLink{sga3:ii:section:1}{Section~1}, there is an isomorphism, functorial in
\(\mathcal M\),
\[
  T_{\underline{\operatorname{Hom}}_{Z/S}(X,Y)/S}(\mathcal M)
  \xrightarrow{\sim}
  \underline{\operatorname{Hom}}_{Z/S}
    \bigl(X,T_{Y/Z}(\mathcal M)\bigr).
\]
This follows immediately from the definitions, in view of the isomorphism
\[
  \underline{\operatorname{Hom}}_S
    \bigl(T,\underline{\operatorname{Hom}}_{Z/S}(X,Y)\bigr)
  \xrightarrow{\sim}
  \underline{\operatorname{Hom}}_{Z/S}
    \bigl(X,\underline{\operatorname{Hom}}_Z(Z\times_ST,Y)\bigr).\footnotemark
\]
\footnotetext{Translator's note: the Gille--Polo re-edition (2011),
Expos\'e~II, spells this out: one takes \(T=I_S(\mathcal M)\), for which the
left-hand side is
\(T_{\underline{\operatorname{Hom}}_{Z/S}(X,Y)/S}(\mathcal M)\), and uses
\(Z\times_SI_S(\mathcal M)\simeq I_Z(\mathcal M)\) (\S2), so that
\(\underline{\operatorname{Hom}}_Z(Z\times_SI_S(\mathcal M),Y)=
T_{Y/Z}(\mathcal M)\).}

\SGAInlineRefTarget{sga3:ii:corollary:3-11-1}\subsection*{Corollary 1.}
If \(Y/Z\) satisfies (E), then
\(\underline{\operatorname{Hom}}_{Z/S}(X,Y)/S\) satisfies (E), and the
isomorphism of \SGARefLink{sga3:ii:proposition:3-11}{3.11} respects the module structures.

\par\medskip\noindent
\emph{Addition} (after the Gille--Polo re-edition (2011), Expos\'e~II; our
formulation). The 1970 text gives no proof of this corollary; here is one.
Put \(H=\underline{\operatorname{Hom}}_{Z/S}(X,Y)\). The functor
\(\underline{\operatorname{Hom}}_{Z/S}(X,-)\) commutes with fibre products,
its value on \(S'\) being \(\operatorname{Hom}_Z(X\times_SS',-)\). If
\(Y/Z\) satisfies (E), it therefore carries
\(T_{Y/Z}(\mathcal M\oplus\mathcal N)\simeq
T_{Y/Z}(\mathcal M)\times_YT_{Y/Z}(\mathcal N)\) to a fibre product over
\(H\), and by 3.11 this says that \(H/S\) satisfies (E). For the module
structures, let \(H'\to H\) correspond to the \(Z\)-morphism
\(f\colon X\times_SH'\to Y\). An \(H\)-morphism
\(H'\to T_{H/S}(\mathcal M)\) is an \(S\)-morphism
\(\varphi\colon I_S(\mathcal M)\times_SH'\to H\) lifting \(H'\to H\), that
is, a \(Z\)-morphism \(\Phi\colon X\times_SI_S(\mathcal M)\times_SH'\to Y\)
extending \(f\), and \(\Phi\) is also the \(Y\)-morphism
\(X\times_SH'\to T_{Y/Z}(\mathcal M)\) that corresponds to \(\varphi\) under
3.11. The operator of \(a\in\mathcal O(H')\) replaces \(\varphi\) by
\(\varphi\circ a^*\) (note to the Corollary of 3.4), that is, \(\Phi\) by
\((x,m,h)\mapsto\Phi(x,m\cdot a(h),h)\). On the other side,
\(b\in\mathcal O(X\times_SH')\) acts through the \(\mathcal O_Y\)-module
\(T_{Y/Z}(\mathcal M)\) (3.6) as in the Addition at the end of \S1, that
is, by \(\Phi\mapsto\bigl((x,m,h)\mapsto\Phi(x,m\cdot b(x,h),h)\bigr)\).
For \(b(x,h)=a(h)\) the two agree, so the isomorphism of 3.11 is
\(\mathcal O_H\)-linear; moreover, the operators of \(\mathcal O(H')\) on
\(\operatorname{Hom}_H(H',T_{H/S}(\mathcal M))\) extend, functorially in
\(H'\), to operators of \(\mathcal O(X\times_SH')\).
\par\medskip

\SGAInlineRefTarget{sga3:ii:corollary:3-11-2}\subsection*{Corollary 2.}
There is an isomorphism, functorial in \(\mathcal M\),
\[
  T_{\underline{\operatorname{Hom}}_S(X,Y)/S}(\mathcal M)
  \xrightarrow{\sim}
  \underline{\operatorname{Hom}}_S(X,T_{Y/S}(\mathcal M)).
\]
If \(Y/S\) satisfies (E), then
\(\underline{\operatorname{Hom}}_S(X,Y)/S\) satisfies (E), and this
isomorphism respects the module structures over
\(\underline{\operatorname{Hom}}_S(X,Y)\).

\SGAInlineRefTarget{sga3:ii:corollary:3-11-3}\subsection*{Corollary 3.}
There is an isomorphism, functorial in \(\mathcal M\),
\[
  L^u_{\underline{\operatorname{Hom}}_S(X,Y)/S}(\mathcal M)
  \xrightarrow{\sim}
  \underline{\operatorname{Hom}}_{Y/S}(X,T_{Y/S}(\mathcal M)),
\]
where \(X\) is regarded as a functor over \(Y\) by means of
\(u:X\to Y\).\footnote{Translator's note: the Gille--Polo re-edition (2011),
Expos\'e~II, deduces this from Corollary~2: \(u\) is a section of
\(\underline{\operatorname{Hom}}_S(X,Y)\) over \(S\), and the fibre over it
of \(\underline{\operatorname{Hom}}_S(X,T_{Y/S}(\mathcal M))\to
\underline{\operatorname{Hom}}_S(X,Y)\) is
\(\underline{\operatorname{Hom}}_{Y/S}(X,T_{Y/S}(\mathcal M))\). It adds that
the isomorphism is \(\mathcal O_S\)-linear when \(Y/S\) satisfies (E), and
likewise that of Corollary~4 (the case \(Y=X\), \(u=\operatorname{id}_X\),
the unit of the monoid \(\underline{\operatorname{End}}_S(X)\)) when \(X/S\)
satisfies (E): by Corollary~2, the isomorphism there is linear over
\(\underline{\operatorname{Hom}}_S(X,Y)\), hence on the fibres over a
section.}

\SGAInlineRefTarget{sga3:ii:corollary:3-11-4}\subsection*{Corollary 4.}
There is an isomorphism, functorial in \(\mathcal M\),
\[
  \underline{\operatorname{Lie}}
    (\underline{\operatorname{End}}_S(X)/S,\mathcal M)
  \xrightarrow{\sim}
  \prod_{X/S}T_{X/S}(\mathcal M).\footnotemark
\]
\footnotetext{Translator's note: the Gille--Polo re-edition (2011),
Expos\'e~II, observes that when \(X/S\) satisfies (E), the right-hand side
\(\prod_{X/S}T_{X/S}(\mathcal M)=
\underline{\operatorname{Hom}}_{X/S}(X,T_{X/S}(\mathcal M))\) is a module
over \(\prod_{X/S}\mathcal O_X\): its value on \(S'\) is a module over
\(\mathcal O(X\times_SS')\), functorially in \(S'\). This follows from 3.6
and the Addition at the end of \S1, or from the Addition after Corollary~1.}

We shall now interpret the definition of the tangent bundle
geometrically. Let \(T\) be a scheme over \(S\); there are isomorphisms,
functorial in \(\mathcal M\),
\[
\begin{aligned}
  T_{X/S}(\mathcal M)(T)
  &\xrightarrow{\sim}\operatorname{Hom}_S\bigl(T,T_{X/S}(\mathcal M)\bigr)
   =\operatorname{Hom}_S\bigl(T,\underline{\operatorname{Hom}}_S(I_S(\mathcal M),X)\bigr)\\
  &\xrightarrow{\sim}\operatorname{Hom}_S\bigl(I_S(\mathcal M),\underline{\operatorname{Hom}}(T,X)\bigr)
   =T_{\underline{\operatorname{Hom}}_S(T,X)/S}(\mathcal M)(S)\\
  &\xrightarrow{\sim}\operatorname{Hom}_{I_S(\mathcal M)}
   \bigl(T_{I_S(\mathcal M)},X_{I_S(\mathcal M)}\bigr).
\end{aligned}
\]
In particular, the morphism \(\mathcal M\to0\) gives a commutative diagram
\[
\begin{tikzcd}[column sep=large]
  \operatorname{Hom}_S\bigl(T,T_{X/S}(\mathcal M)\bigr) \arrow[r,"\sim"] \arrow[d]
  & \operatorname{Hom}_{I_S(\mathcal M)}\bigl(T_{I_S(\mathcal M)},X_{I_S(\mathcal M)}\bigr) \arrow[d]\\
  \operatorname{Hom}_S(T,X) \arrow[r,"\text{identity}"]
  & \operatorname{Hom}_S(T,X)
\end{tikzcd}
\]
where the first vertical arrow is obtained from the projection
\(T_{X/S}(\mathcal M)\to X\), and the second from the base change
\(S\to I_S(\mathcal M)\). Consequently:\footnote{Translator's note: the
Gille--Polo re-edition (2011), Expos\'e~II, takes for \(T\) an arbitrary
\(S\)-functor \(U\), the chain of isomorphisms becoming
\(\operatorname{Hom}_S(U,\underline{\operatorname{Hom}}_S(I_S(\mathcal M),X))
\simeq\operatorname{Hom}_S(I_S(\mathcal M),\underline{\operatorname{Hom}}_S(U,X))
\simeq\operatorname{Hom}_{I_S(\mathcal M)}(U_{I_S(\mathcal M)},X_{I_S(\mathcal M)})\)
by \SGARefLink{sga3:i:corollary:1-7-2}{I~1.7.2}. This covers the corollary below, which applies the
proposition with \(T=X\), a functor that need not be representable.}

\SGAInlineRefTarget{sga3:ii:proposition:3-12}\subsection*{\SGARefLink{sga3:ii:proposition:3-12}{Proposition 3.12}.}
Let \(T\) be over \(X\) by \(h_0:T\to X\). The set
\(\operatorname{Hom}_X(T,T_{X/S}(\mathcal M))\) identifies with the set of
\(I_S(\mathcal M)\)-morphisms from \(T_{I_S(\mathcal M)}\) to
\(X_{I_S(\mathcal M)}\) that restrict to \(h_0\) over
\(S\subset I_S(\mathcal M)\).

\SGAInlineRefTarget{sga3:ii:corollary:3-12-1}\subsection*{Corollary.}
The set \(\Gamma(T_{X/S}(\mathcal M)/X)\) identifies with the set of
\(I_S(\mathcal M)\)-endomorphisms of \(X_{I_S(\mathcal M)}\) that induce
the identity on \(X\).

Let us now observe that, by the preceding isomorphisms, this last set also
identifies with
\(\underline{\operatorname{Lie}}(\underline{\operatorname{End}}_S(X)/S,\mathcal M)(S)\).
If \(X/S\) satisfies (E), then \(\underline{\operatorname{End}}_S(X)/S\)
satisfies (E) (\SGARefLink{sga3:ii:corollary:3-11-2}{3.11, Corollary~2}). Applying
\SGARefLink{sga3:ii:corollary:3-9-2}{3.10, Corollary~2}, and \SGARefLink{sga3:ii:proposition:3-9}{3.9}, one deduces the following.

\SGAInlineRefTarget{sga3:ii:proposition:3-13}\subsection*{\SGARefLink{sga3:ii:proposition:3-13}{Proposition 3.13}.}
If \(X/S\) satisfies (E), every \(I_S(\mathcal M)\)-endomorphism of
\(X_{I_S(\mathcal M)}\) that induces the identity on \(X\) is an
automorphism. The group \(\Gamma(T_{X/S}(\mathcal M)/X)\) identifies with
the group of these automorphisms.

\SGAInlineRefTarget{sga3:ii:corollary:3-13-1}\subsection*{Corollary 1.}
Let \(u:X\to Y\) be an \(S\)-isomorphism, with \(Y/S\) satisfying (E).
Every \(I_S(\mathcal M)\)-morphism
\(X_{I_S(\mathcal M)}\to Y_{I_S(\mathcal M)}\) extending \(u\) is an
isomorphism.

\SGAInlineRefTarget{sga3:ii:corollary:3-13-2}\subsection*{Corollary 2.}
Let \(u\in\underline{\operatorname{Isom}}_S(X,Y)\), with \(Y/S\)
satisfying (E). The monomorphism
\[
  \underline{\operatorname{Isom}}_S(X,Y)
  \longrightarrow\underline{\operatorname{Hom}}_S(X,Y)
\]
induces an isomorphism
\[
  L^u_{\underline{\operatorname{Isom}}_S(X,Y)/S}(\mathcal M)
  \xrightarrow{\sim}
  L^u_{\underline{\operatorname{Hom}}_S(X,Y)/S}(\mathcal M).\footnotemark
\]
\footnotetext{Translator's note: the Gille--Polo re-edition (2011),
Expos\'e~II, adds a proof. In our words: by 3.4 it suffices to compare the
points with values in \(S\), after any base change \(S'\to S\); by
Corollary~3 of 3.11 and by 3.12, those on the right are the
\(I_S(\mathcal M)\)-morphisms \(X_{I_S(\mathcal M)}\to Y_{I_S(\mathcal M)}\)
extending \(u\), and those on the left are the ones among them that are
isomorphisms; by Corollary~1, all of them are.}

\SGAInlineRefTarget{sga3:ii:corollary:3-13-3}\subsection*{Corollary 3.}
Let \(u\in\underline{\operatorname{Aut}}_S(X)\), with \(X/S\) satisfying
(E). The monomorphism
\(\underline{\operatorname{Aut}}_S(X)\to
\underline{\operatorname{End}}_S(X)\) induces an isomorphism
\[
  L^u_{\underline{\operatorname{Aut}}_S(X)/S}(\mathcal M)
  \xrightarrow{\sim}
  L^u_{\underline{\operatorname{End}}_S(X)/S}(\mathcal M).
\]
In particular, one has
\[
  \underline{\operatorname{Lie}}
    (\underline{\operatorname{Aut}}_S(X)/S,\mathcal M)(S)
  \xrightarrow{\sim}
  \underline{\operatorname{Lie}}
    (\underline{\operatorname{End}}_S(X)/S,\mathcal M)(S)
  \xrightarrow{\sim}
  \prod_{X/S}\bigl(T_{X/S}(\mathcal M)\bigr).\footnotemark
\]
\footnotetext{Translator's note: the Gille--Polo re-edition (2011),
Expos\'e~II, states these isomorphisms for the functors themselves, which
follows from the isomorphisms of the corollary (isomorphisms of
\(S\)-functors) and from Corollary~4 of 3.11, and adds that
\(\underline{\operatorname{Lie}}(\underline{\operatorname{Aut}}_S(X)/S,\mathcal M)\)
thereby becomes a module over \(\prod_{X/S}\mathcal O_X\) when \(X/S\)
satisfies (E) (see our note to Corollary~4 of 3.11).}

\subsection*{\SGAInlineRefTarget{sga3:ii:subsection:3-14}3.14.}
Finally, suppose that \(X\) is representable. We have seen that a
section of \(T_{X/S}\) over \(X\) identifies with an
\(I_S\)-automorphism of \(X_{I_S}\) inducing the identity on \(X\). Now
\(X_{I_S}\) and \(X\) have the same underlying topological space, the
corresponding sheaves of rings being \(\mathscr D_{\mathcal O_X}\) and
\(\mathcal O_X\). It follows at once that an infinitesimal automorphism of
\(X\) identifies with a derivation of the sheaf of rings \(\mathcal O_X\)
over the sheaf \(\mathcal O_S\). We have thus constructed an isomorphism
\[
  \Gamma(T_{X/S}/X)
  \xrightarrow{\sim}\operatorname{Der}_{\mathcal O_S}(\mathcal O_X)
\]
which, by \SGARefLink{sga3:ii:proposition:3-13}{3.13}, respects the group structures (and even the
\(\mathcal O(S)\)-module structures), recovering the classical
interpretation of tangent vector fields as derivations of the structure
sheaf.\footnote{Translator's note: the Gille--Polo re-edition (2011),
Expos\'e~II, expands this paragraph. In our words: by the Nota of 2.2,
\(T_{X/S}=V(\Omega^1_{X/S})\), so
\(\Gamma(T_{X/S}/X)=\operatorname{Hom}_{\mathcal O_X}(\Omega^1_{X/S},\mathcal O_X)
=\operatorname{Der}_{\mathcal O_S}(\mathcal O_X)\) directly. Along the route
above, an endomorphism of \(X_{I_S}\) over \(I_S\) inducing the identity on
\(X\) is given by a morphism of \(\mathcal O_S\)-algebras
\(\mathcal O_X\to\mathcal O_X\oplus\varepsilon\mathcal O_X\) lifting the
identity, that is, \(f\mapsto f+\varepsilon D(f)\) with \(D\) an
\(\mathcal O_S\)-derivation; extended \(\varepsilon\)-linearly, the
endomorphism attached to \(D+D'\) is the composite of those attached to
\(D\) and \(D'\), and the operator of \(a\in\mathcal O(X)\) turns \(D\) into
\(aD\), so the bijection is even \(\mathcal O(X)\)-linear. The re-edition
also notes its compatibility with base change: for \(S'\to S\) and
\(X'=X_{S'}\),
\(\underline{\operatorname{Lie}}(\underline{\operatorname{Aut}}_S(X)/S)(S')
\simeq\Gamma(T_{X'/S'}/X')\simeq\operatorname{Der}_{\mathcal O_{S'}}(\mathcal O_{X'})\).}
Let us note, moreover, that
\[
  \Gamma(T_{X/S}/X)=H^0(X,\mathfrak g_{X/S}),
\]
where \(\mathfrak g_{X/S}\) is the dual of \(\Omega^1_{X/S}\).

